% Part: computability
% Chapter: computability-theory
% Section: total

\documentclass[../../../include/open-logic-section]{subfiles}

\begin{document}

\olfileid[hi]{cmp}{thy}{tot}
\olsection{पूर्णता अनिर्णेय है}

किसी वस्तु को असंगणनीय दिखाने के लिए $s$-$m$-$n$ प्रमेय के उपयोग का एक
और उदाहरण देखें। $\fn{Tot}$ को पूर्ण संगणनीय फलनों के सूचकांकों का
समुच्चय मानें, अर्थात्
\[
\fn{Tot} = \Setabs{x}{\text{प्रत्येक $y$ के लिए, $\cfind{x}(y)\fdefined$}}.
\]

\begin{prop}
\ollabel{prop:total}
$\fn{Tot}$ संगणनीय नहीं है।
\end{prop}

\begin{proof}
$\fn{Tot}$ को असंगणनीय देखने के लिए यह दिखाना पर्याप्त है कि $K$ उसमें
अपचयनीय है। $h(x,y)$ को परिभाषित करें:
\[
h(x,y) \simeq
\begin{cases}
0 & \text{यदि $x \in K$} \\
\fundefined & \text{अन्यथा}
\end{cases}
\]
ध्यान दें कि $h(x,y)$ तनिक भी $y$ पर निर्भर नहीं है। यह देखना कठिन नहीं
कि $h$~आंशिक संगणनीय है: आगत $x, y$ पर~$h$ को संगणित करने के लिए पहले
फलन~$\cfind{x}$ का आगत~$x$ पर अनुकरण करते हैं; यदि यह संगणना रुकती है,
तो $h(x,y)$, $0$ निर्गत करके रुक जाता है। अतः $h(x,y)$ केवल
$\Zero(\umin{s}{T(x,x,s)})$ है, जहाँ $\Zero$ अचर शून्य फलन है।

$s$-$m$-$n$ प्रमेय से कोई आदिम पुनरावर्ती फलन~$k(x)$ ऐसा है कि प्रत्येक
$x$ और~$y$ के लिए
\[
\cfind{k(x)}(y) =
\begin{cases}
0 & \text{यदि $x \in K$} \\
\fundefined & \text{अन्यथा}
\end{cases}
\]
अतः $\cfind{k(x)}$ पूर्ण है यदि $x \in K$, और अन्यथा अपरिभाषित।
इसलिए $k$, $K$ से~$\fn{Tot}$ में एक अपचयन है।
\end{proof}

\begin{digress}
वास्तव में $\fn{Tot}$ संगणनीयतः परिगणनीय तक नहीं है---उसकी जटिलता
``अंकगणितीय पदानुक्रम'' में और ऊपर स्थित है। किंतु यहाँ हम इस प्रबलतर
परिणाम की चिंता नहीं करेंगे।
\end{digress}

\end{document}

